\section{Introduction}
\label{sec:introduction}

The Yang-Mills existence and mass gap problem, one of the seven Millennium Prize Problems, asks for a rigorous mathematical proof that non-Abelian gauge theories in four dimensions possess a mass gap—a strictly positive lower bound on the energy spectrum above the vacuum state.

\subsection{Scope of this Work}

This manuscript presents a rigorous framework for analyzing the problem, clearly distinguishing between established results in constructive quantum field theory and the remaining open problems required for a full solution.
We do not assert an unconditional construction of continuum four-dimensional Yang--Mills theory in this version. Instead, we aim to:
\begin{enumerate}
    \item Establish the rigorous lattice foundation (Section \ref{sec:lattice}).
    \item Prove the existence of the mass gap in the strong coupling regime via cluster expansion (Section \ref{sec:strong_coupling}).
    \item Formulate precise conditional strategies for bridging the gap to weak coupling and the continuum limit (Section \ref{sec:adjoint}).
\end{enumerate}

\subsection{Architecture of the Analysis}

The analysis is organized into distinct regimes:

\subsubsection*{Stage I: The Lattice Foundation (Finite Volume)}
We begin with the standard Wilson lattice regularization. Using the transfer matrix formalism and reflection positivity, we review the existence of a spectral gap $\Delta_L(\beta) > 0$ for any finite lattice volume $L^3$ and any coupling $\beta$.

\subsubsection*{Stage II: Strong Coupling Control ($\beta < \beta_c$)}
In the strong coupling regime, we utilize the \textbf{Cluster Expansion} technique. We provide a rigorous derivation of the convergence of the polymer expansion of the partition function. This establishes the exponential decay of correlations and a mass gap $\Delta \propto \ln(1/\beta)$ in this regime.

\subsubsection*{Stage III: The Challenge of Intermediate Coupling}
To extend these results to the continuum limit ($\beta \to \infty$), one must control the intermediate coupling regime. We discuss the \textbf{Adjoint Interpolation} strategy, which proposes to connect Pure Yang-Mills theory to Supersymmetric Yang-Mills theory (where a gap is known). We state the precise analytic hypotheses required for this strategy to succeed, specifically the absence of phase transitions along the interpolation path.

\subsubsection*{Stage IV: The Continuum Limit}
We discuss the formal requirements for the continuum limit, including the reconstruction of the Hilbert space via Osterwalder-Schrader reconstruction and the necessary uniform bounds on the renormalization group flow.

\subsection{Conventions}
Throughout this paper, we consider $G=SU(N)$ gauge theory on a four-dimensional Euclidean lattice $\Lambda$. We use natural units $\hbar = c = 1$.
